% GENERATED FILE. Edit canonical lessons, then run npm run generate:books.
% canonical-source-hash: fnv1a64:4ab358f1746925ec
% canonical-lessons: ES-C09-estamos, ES-C09-estan, ES-C09-estais, ES-C09-repaso-ser-estar

\chapter{Estar in the Plural}
\label{ch:spanish-estar-plural}
\hlchaptermodality{50}

\begin{chapteropening}
\textbf{By the end of this chapter:} \emph{I can use all six present forms of estar, and say why its plural is regular where ser's is not.}

One Latin verb against two - which is the whole reason one paradigm is plain and the other is not.
\end{chapteropening}

\section[estamos]{estamos — the other \textquotedblleft{}we are\textquotedblright{}}

\label{lesson:ES-C09-estamos}



\begin{quote}
Say \textquotedblleft{}we are\textquotedblright{} with \emph{ser}. (\emph{Somos}.) And how do you ask a friend how they are? (\emph{¿Cómo estás?})
\end{quote}

\begin{grammarlens}[title={Grammar Lens: this one behaves}]
\begin{quote}
\textbf{estamos} — \emph{we are.}
\end{quote}

Drop the \textbf{-ar}, add \textbf{-amos}. That is the plain \emph{-ar} ending you learned on \emph{hablamos}, applied without alteration.

And that is the point of this lesson. You have just spent a chapter on \emph{ser}, whose six forms come from two different Latin verbs and refuse to resemble each other. \textbf{Estar is not like that at all.} In the plural it is an ordinary \emph{-ar} verb.

\begin{quote}
\emph{hablamos} · \emph{estamos} — the same ending, and no surprises.
\end{quote}
\end{grammarlens}

\begin{cousinweb}
The reason is simple, and it is the reason for almost every irregularity you have met: \textbf{estar is one verb.} Latin \emph{stāre}, \textquotedblleft{}to stand\textquotedblright{} — and that is the whole ancestry. Nothing was fused into it, so nothing in it has a seam.

\emph{Ser} is irregular because two verbs were welded together. \emph{Estar} is regular because it never was.

That is worth carrying, because it predicts the rest. Where a Spanish verb looks chaotic, the usual explanation is not that it decayed strangely but that it is \textbf{more than one word wearing a single name}.
\end{cousinweb}

\subsection*{Your turn}
Say these aloud:

\begin{itemize}
\raggedright
  \item \textquotedblleft{}hablamos, estamos\textquotedblright{} — the same ending
  \item \textquotedblleft{}somos, estamos\textquotedblright{} — two verbs for \emph{we are}, and only one of them is strange
\end{itemize}

\begin{tcolorbox}[breakable,colback=teal!4,colframe=teal!35!black,title={Before you move on}]
What is \textquotedblleft{}we are\textquotedblright{} with \emph{estar}? (\textbf{Estamos}.) Which Latin verb is it from? (\emph{\textbf{Stāre}}, \textquotedblleft{}to stand\textquotedblright{}.) And why is it regular where \emph{ser} is not? (Because it is \textbf{one} verb, not two.)
\end{tcolorbox}
\section[están]{están — \textquotedblleft{}they are\textquotedblright{}}

\label{lesson:ES-C09-estan}



\begin{quote}
Say \textquotedblleft{}we are\textquotedblright{} with \emph{estar}. (\emph{Estamos}.) And \textquotedblleft{}they speak.\textquotedblright{} (\emph{Hablan}.)
\end{quote}

\begin{grammarlens}[title={Grammar Lens: the ending is regular, the accent is the clue}]
\begin{quote}
\textbf{están} — \emph{they are}, and \textbf{ustedes están}, \emph{you all are}.
\end{quote}

The ending is \textbf{-an}, exactly as in \emph{hablan}. But there is a written accent, and it is not decoration.

\emph{Estar} is unusual in one specific way: it puts its stress on the \textbf{ending} rather than the stem, in every singular and plural form except \emph{estamos}. You have already seen this without being told — \emph{estás}, \emph{está}, and now \emph{están}. The accent is simply Spanish writing that stress down, by the rule for written accents you learned long ago.

\begin{quote}
\textbf{hablan} — stress on the stem, no accent. \textbf{están} — stress on the ending, so an accent.
\end{quote}

Same ending, different stress, and the spelling records it faithfully.
\end{grammarlens}

\subsection*{Your turn}
Say these aloud:

\begin{itemize}
\raggedright
  \item \textquotedblleft{}hablan\textquotedblright{} then \textquotedblleft{}están\textquotedblright{} — same ending, and hear the stress move
  \item \textquotedblleft{}ustedes están\textquotedblright{}, addressing a group
\end{itemize}

\begin{tcolorbox}[breakable,colback=teal!4,colframe=teal!35!black,title={Before you move on}]
What is \textquotedblleft{}they are\textquotedblright{} with \emph{estar}? (\textbf{Están}.) Why the accent? (The stress falls on the \textbf{ending}.) Which form of \emph{estar} does not do that? (\emph{\textbf{Estamos}}.)
\end{tcolorbox}
\section[estáis]{estáis — and nothing surprises you}

\label{lesson:ES-C09-estais}



\begin{quote}
Say \textquotedblleft{}they are\textquotedblright{} with \emph{estar}. (\emph{Están}.) And the \emph{vosotros} form of \emph{hablar}? (\emph{Habláis}.)
\end{quote}

\begin{grammarlens}[title={Grammar Lens: predictable, and that is the news}]
\begin{quote}
\textbf{vosotros estáis} — \emph{you all are}, in Spain.
\end{quote}

If you guessed that before reading it, you have understood the chapter. The ending is \emph{-áis}, the plain \emph{-ar} one from \emph{habláis}, and the accent is there for the reason it is always there with this verb: the stress sits on the ending.

That is now three forms in a row where knowing \emph{estar} meant knowing nothing new. In the Americas, \emph{ustedes están} does this job.
\end{grammarlens}

\subsection*{Your turn}
\begin{itemize}
\raggedright
  \item \emph{Say it:} \textquotedblleft{}habláis\textquotedblright{} then \textquotedblleft{}estáis\textquotedblright{} — the same ending, twice
\end{itemize}

\begin{tcolorbox}[breakable,colback=teal!4,colframe=teal!35!black,title={Before you move on}]
What is the \emph{vosotros} form of \emph{estar}? (\textbf{Estáis}.) Why the accent? (\emph{Estar} stresses its \textbf{ending}.) And in Argentina? (\emph{\textbf{Ustedes están}}.)
\end{tcolorbox}
\section[Practice]{Review — both verbs for \textquotedblleft{}to be\textquotedblright{}}

\label{lesson:ES-C09-repaso-ser-estar}



\begin{quote}
Nothing new. Say \emph{somos}, then \emph{estamos}.
\end{quote}

\begin{grammarlens}[title={Grammar Lens: side by side, and the difference is history}]
\noindent
\begin{tabularx}{\linewidth}{@{}>{\raggedright\arraybackslash}X>{\raggedright\arraybackslash}X>{\raggedright\arraybackslash}X@{}}
\toprule
\textbf{who} & \textbf{\emph{ser}} & \textbf{\emph{estar}} \\
\midrule
yo & \textbf{soy} & \textbf{estoy} \\
tú & \textbf{eres} & \textbf{estás} \\
él · ella · usted & \textbf{es} & \textbf{está} \\
nosotros & \textbf{somos} & \textbf{estamos} \\
vosotros & \textbf{sois} & \textbf{estáis} \\
ellos · ustedes & \textbf{son} & \textbf{están} \\
\bottomrule
\end{tabularx}

Read the two columns down and the difference is not a matter of difficulty. It is a matter of \textbf{history}.

\begin{itemize}
\raggedright
  \item \emph{\textbf{Ser}} is two Latin verbs fused — \emph{esse} and \emph{sedēre} — so its forms have no common shape. Five begin with \emph{s-}; \emph{eres} does not, and that is the seam.
  \item \emph{\textbf{Estar}} is one Latin verb, \emph{stāre}. Its plural is a plain \emph{-ar} paradigm, and the only oddity in the whole column is where the stress falls, which the written accent records honestly.
\end{itemize}

So the two hardest-looking verbs in Spanish are hard for opposite reasons: one because it is two words, and the other not really at all.
\end{grammarlens}

\subsection*{Your turn}
Say these aloud:

\begin{itemize}
\raggedright
  \item one column down, then the other
  \item each row across — \emph{soy, estoy} · \emph{somos, estamos} — and hear which column surprises you
\end{itemize}

\emph{Twice through:}

\begin{tcolorbox}[breakable,colback=teal!4,colframe=teal!35!black,title={Before you move on}]
Which of the two verbs is regular in the plural? (\emph{\textbf{Estar}}.) Why? (It is \textbf{one} Latin verb; \emph{ser} is two.) And what is the only oddity in \emph{estar}'s column? (Where the \textbf{stress} falls — which the accent records.)
\end{tcolorbox}
